\section{PCA and Dimension-Ordering Analysis}
\label{sec:results-pca}

If late pruning is caused by energy being spread across dimensions, then a
rotation that concentrates energy in the leading dimensions should make the
BOND bound tighten earlier. This section tests that mechanism hypothesis with
two diagnostic interventions. Both are explicitly \emph{diagnostic}: the
oracle seed policy is not deployable, and while PCA itself is a legitimate
offline preprocessing step, here it is used to isolate the effect of dimension
ordering on the bound.

\subsection{Oracle seeds}

The implementable threshold uses the first 500 corpus documents as
deterministic seeds, so seed quality could in principle limit pruning.
To remove this confounder, an oracle arm receives the exact top-10 documents
as \emph{free} seeds---their scoring cost is not charged---which yields the
tightest threshold any seed policy could produce. On raw held-out SciFact
embeddings (\cref{tab:bond}, second row), oracle seeding raises document
pruning from 22.44\% to 36.32\%, but the evaluated component products drop
only from 95.47\% to 92.43\%, because every prune still occurs at checkpoint
96. The oracle arm is in fact marginally \emph{slower}
(\SI{45.09}{\second} median, $11.50\times$ the same-run exhaustive arm):
earlier threshold tightening increases bound-check activity without saving
meaningful arithmetic. Seed selection is therefore not the bottleneck; the
bound geometry is.

\subsection{Offline PCA rotation}

An uncentered PCA is fitted offline on the document token vectors and the
resulting orthogonal rotation is applied offline to both documents and
queries. Because the rotation is orthogonal, it preserves inner products in
exact arithmetic (measured float32 orthogonality error
$2.53 \times 10^{-8}$), so the exact ranking is unchanged. The offline costs
are small and reported separately: \SI{0.592}{\second} to fit,
\SI{0.227}{\second} to transform, and \SI{1.692}{\second} to build the
vertical index. After rotation, the first 8 of 128 components carry 90.91\%
of total token energy, and the first 64 carry 97.76\%.

The effect on the mechanism is substantial (\cref{tab:bond}, rows 3--4, and
\cref{fig:bond}). With the implementable prefix-500 seeds, document pruning
rises to 86.30\% and evaluated component products fall to 69.02\%; with free
oracle seeds, 97.91\% of documents prune and only 61.53\% of products are
evaluated. Pruning also genuinely moves earlier: most prunes now occur at
checkpoints 64 and 96 rather than exclusively at 96.

\begin{figure}[tb]
  \centering
  \includegraphics[width=\textwidth]{fig7_controlled_bond.png}
  \caption{Exact-safe BOND-MaxSim on 40 held-out SciFact queries (5{,}183
    documents; four pinned physical cores; median of five interleaved
    repetitions; oracle seeds are free and diagnostic). Left: fraction of
    exhaustive component products evaluated (solid) and query--document pairs
    pruned (hatched) for raw and PCA-rotated embeddings under prefix-500 and
    oracle seed policies. Right: BOND median latency relative to the
    precision-matched exhaustive median of the same run. PCA rotation improves
    pruning substantially (69.0\% and 61.5\% of products evaluated) yet every
    arm remains $8.4$--$11.5\times$ slower than exhaustive scoring.}
  \label{fig:bond}
\end{figure}

\subsection{Arithmetic reduction does not become latency}

Despite removing 31--38\% of the component products, both PCA arms remain
8.38--8.51 times slower than the exhaustive arm interleaved in the same run
(\SI{36.95}{\second} and \SI{37.55}{\second} versus \SI{4.41}{\second}
median). The latency gap narrows from $10.76\times$ to 8.38--8.51 times while the
arithmetic drops by nearly a third, which localizes the remaining cost: bound
maintenance, threshold updates, branching, partial-state reads and writes, and
the irregular access pattern of surviving documents dominate the runtime of
this kernel, not the multiply--accumulate work itself. Dimension ordering
demonstrably helps the \emph{bound}; it does not make this \emph{kernel}
competitive. This completes the answer to RQ2 and explains why inverse
operation counts (as reported by the earlier NumPy instrumentation) must not
be read as speedups.
